\NSPage{091}%
\end{lemma}
\begin{proof}
We first estimate the primitive without differentiating its rapidly shifted argument.
We then verify the exact radial identity and use torus Fourier inversion to estimate its cutoff
remainder.

\noindent\textbf{Step 1: support and coefficient estimates.} Work on a neighborhood with a fixed common
index, as in Lemma~\ref{lem:6.2}, and put \NSi{NS-F-P091-001}{\mathcal L=v_r\cdot\partial_y} and
\NSi{NS-F-P091-002}{g=R^e f}. Multiplication by \NSi{NS-F-P091-003}{R^e} preserves
the input estimate on the shell. To estimate \NSi{NS-F-P091-004}{\mathcal I_cg}, set
\NSi{NS-F-P091-005}{U=R^{d_r}} and smoothly extend
\NSi{NS-F-P091-006}{F(U)=g(R)/(d_rR^{d_r-1})} by zero. The shell stays in a fixed compact subset
of \NSi{NS-F-P091-007}{R>0}, so this change of variables and its inverse have bounded derivatives
of every fixed order. The exact half-line representation is
\NSFormula{NS-F-P091-008}{display}{8.9}
\begin{equation}
\begin{aligned}
A_-(U,y)&=\int_{-\infty}^{0}F(U+u,y+Muv_r)\,du,\\
A_+(U,y)&=\int_{0}^{\infty}F(U+u,y+Muv_r)\,du,\\
\mathcal Jg&=A_-+A_+,\qquad
\mathcal I_cg=(1-\chi_m)A_- -\chi_mA_+.
\end{aligned}
\tag{8.9}\label{eq:8.9}
\end{equation}
At fixed \NSi{NS-F-P091-009}{u}, the shift depends on neither \NSi{NS-F-P091-010}{U} nor
\NSi{NS-F-P091-011}{(Z,T)}. Consequently a coefficient derivative differentiates only
\NSi{NS-F-P091-012}{F} and \NSi{NS-F-P091-013}{\chi_m}, without a factor
\NSi{NS-F-P091-014}{M}. The input zero extension includes derivatives of the moving shell, so
differentiation creates no boundary terms.

Near the inner edge only the left integral occurs. Write
\NSi{NS-F-P091-015}{d_a=\log(X/X_a)}, \NSi{NS-F-P091-016}{d_b=\log(X_b/X)}, so
\NSi{NS-F-P091-017}{\zeta=\exp(-a_a/d_a^2-a_b/d_b^2)} and
\NSi{NS-F-P091-018}{\delta=\min(1,d_a,d_b)}, with fixed positive
\NSi{NS-F-P091-019}{a_a,a_b}. For every fixed \NSi{NS-F-P091-020}{B}, the function
\NSi{NS-F-P091-021}{s\mapsto e^{-a_a/s^2}s^{-B}} is increasing for sufficiently small positive
\NSi{NS-F-P091-022}{s}. On the left integration segment this bounds the input weight by the same
type of weight at its endpoint; the factors belonging to the opposite edge are bounded and
comparable. Near the outer edge use the right integral and the analogous monotonicity for
\NSi{NS-F-P091-023}{d_b}. In the interior the output weight is bounded below. Integration takes
place over bounded length. These observations apply after every fixed coefficient derivative,
proving the \NSi{NS-F-P091-024}{\mathsf M_\alpha} estimate. Formula~\eqref{eq:8.9} also makes the
output zero below and above the shell. Multiplication by \NSi{NS-F-P091-025}{R^{\pm e}} costs only
fixed constants there.

\noindent\textbf{Step 2: the radial equation.} The operator
\NSi{NS-F-P091-026}{\partial_U+M\mathcal L} on the integrand in \eqref{eq:8.9} equals its
\NSi{NS-F-P091-027}{u}-derivative. Thus it sends \NSi{NS-F-P091-028}{A_-} to
\NSi{NS-F-P091-029}{F}, \NSi{NS-F-P091-030}{A_+} to \NSi{NS-F-P091-031}{-F}, and
\NSi{NS-F-P091-032}{\mathcal Jg} to zero. Multiplication by
\NSi{NS-F-P091-033}{d_rR^{d_r-1}}, followed by the product rule for
\NSi{NS-F-P091-034}{R^{-e}}, gives \eqref{eq:8.7}. In physical variables the same computation
reads \NSi{NS-F-P091-035}{\mathfrak r\mathcal I_cg=g-(\partial_r\chi_m)\mathcal Jg}.

\noindent\textbf{Step 3: the cutoff remainder under the moment condition.} The radial identity and
the support bounds are now established. To estimate the remaining cutoff term, we use cancellation
in the full radial integral. Haar measure is translation invariant, and hence
\NSi{NS-F-P091-036}{\langle\mathcal Jg\rangle_Y=\int\langle g\rangle_Y\,dR}. Under the weighted
moment condition this is zero for \NSi{NS-F-P091-037}{g=R^ef}. For a zero-mean torus function
\NSi{NS-F-P091-038}{H}, the Diophantine inequality \eqref{eq:6.7} gives
\NSFormula{NS-F-P091-039}{display}{8.10}
\begin{equation}
\|\mathcal L^{-p}H\|_{C_y^m}\le C_{m,p}\|H\|_{C_y^{m+p+3}},
\qquad p\ge1.
\tag{8.10}\label{eq:8.10}
\end{equation}
Indeed the Fourier multiplier is \NSi{NS-F-P091-040}{(2\pi i v_r\cdot k)^{-p}}; after estimating
\NSi{NS-F-P091-041}{m+p+3} derivatives of \NSi{NS-F-P091-042}{H}, the remaining series is bounded
by \NSi{NS-F-P091-043}{\sum_{k\in\mathbb Z^2}(1+|k|)^{-3}}. Slow and radial derivatives commute
with this multiplier. The zero-average part of the full-line integral therefore satisfies the exact
identity
\NSFormula{NS-F-P091-044}{display}{none}
\[
\mathcal Jg=(-M^{-1})^p\int_{\mathbb R}\partial_U^p\mathcal L^{-p}F^\circ
   (U+u,y+Muv_r)\,du.
\]
It follows by integrating the \NSi{NS-F-P091-045}{u}-derivative of
\NSi{NS-F-P091-046}{\mathcal L^{-1}F^\circ} and repeating; every endpoint term vanishes by compact
support. The omitted zero Fourier mode has identically zero full-line
\NSPage{092}%
integral at every slow point. Averaging, integration, and slow differentiation of coefficients
commute on the smooth zero extensions, so all its derivatives have zero integral as well. The
same fixed-shift argument handles further coefficient derivatives. For the output derivative
indexed by \NSi{NS-F-P092-001}{I}, radial input derivatives up to
\NSi{NS-F-P092-002}{I_R+p} and torus derivatives up to
\NSi{NS-F-P092-003}{|I_y|+p+3} suffice, together with the slow derivatives specified by
\NSi{NS-F-P092-004}{I}. Since
\NSi{NS-F-P092-005}{M^{-1}\le C\varepsilon^{\kappa_s}S_*^{\rho_g}}, and
\NSi{NS-F-P092-006}{\partial_R\chi_m} has interior support, this proves \eqref{eq:8.8}. Fix a
physical normalization and a physical derivative order, and let \NSi{NS-F-P092-007}{L\ge0}
bound the resulting loss of powers of \NSi{NS-F-P092-008}{q}, including differentiation after
phase evaluation. The graph operators give such a finite \NSi{NS-F-P092-009}{L}, independent of
\NSi{NS-F-P092-010}{p}. For any prescribed flatness order \NSi{NS-F-P092-011}{N\ge0}, choose
\NSi{NS-F-P092-012}{p} with
\NSi{NS-F-P092-013}{h(\alpha+p\kappa_s)>N+L}. Since
\NSi{NS-F-P092-014}{\varepsilon=Q^h} and \NSi{NS-F-P092-015}{q\asymp Q}, the strict excess
absorbs the fixed powers of \NSi{NS-F-P092-016}{S_*} and gives an
\NSi{NS-F-P092-017}{O(q^N)} bound for that physical derivative. No estimate uniform in
\NSi{NS-F-P092-018}{p} is asserted. Finally, if \NSi{NS-F-P092-019}{f} is independent of
\NSi{NS-F-P092-020}{Y}, then
\NSi{NS-F-P092-021}{\mathcal J(R^ef)=\int R^ef\,dR}, which vanishes under the stated weighted
integral condition.
\end{proof}

The inverse is therefore exact up to a flat cutoff remainder whenever the weighted mean integral
vanishes. This flatness uses estimates at every fixed derivative order. Frequencies comparable to
\NSi{NS-F-P092-022}{M} can be nearly resonant with \NSi{NS-F-P092-023}{v_r}; their small Fourier
coefficients are controlled by the additional derivatives in the torus variables in \eqref{eq:8.10}.
A fixed finite regularity bound would give only a fixed finite flatness order.

\subsection{Pressure reconstruction and divergence-free velocity corrections.}\label{sec:8.3}
In this subsection we will apply the radial inverse of Lemma~\ref{lem:8.2} first to pressure
\NSi{NS-F-P092-024}{p_m}, then to an azimuthal vector potential that produces a divergence-free
velocity increment (Proposition~\ref{prop:8.3}). Its axial component will equal the prescribed
source \NSi{NS-F-P092-025}{\gamma_d} up to the cutoff remainder in \eqref{eq:8.14}. The pressure
source may have a nonzero radial integral of its auxiliary mean. Subtracting a fixed bump times that
integral isolates the obstruction and allows the remaining source to be integrated with compact
support, as in \eqref{eq:8.12}.

Let \NSi{NS-F-P092-026}{\mathcal I_{\mathrm{mean}}} be the reserved interval in \NSi{NS-F-P092-027}{X}
from Theorem~\ref{thm:4.6}(vi), and set
\NSFormula{NS-F-P092-028}{display}{none}
\[
I_m=\{\sqrt{2X}:X\in \mathcal I_{\mathrm{mean}}\}.
\]
The \emph{mean patch} is the region \NSi{NS-F-P092-029}{r/\sqrt q\in I_m}. Choose a fixed bump
\NSi{NS-F-P092-030}{\widehat\rho\in C_c^\infty(I_m)}, normalized by
\NSi{NS-F-P092-031}{\int\widehat\rho(x)\,dx=1}. Define
\NSFormula{NS-F-P092-032}{display}{none}
\[
\rho_{\mathrm{phys}}(r,z,t)=q^{-1/2}\widehat\rho(r/\sqrt q),
\qquad \rho=Q^{1/2}\rho_{\mathrm{phys}}.
\]
The bump is independent of the angular and auxiliary variables, and
\NSi{NS-F-P092-033}{\int\rho_{\mathrm{phys}}\,dr=\int\rho\,dR=1}. It is fixed independently of
the velocity corrections.

The operator \NSi{NS-F-P092-034}{\mathsf T_0} reconstructs pressure, and
\NSi{NS-F-P092-035}{\mathsf T_1} reconstructs an azimuthal vector potential. The operators
\NSi{NS-F-P092-036}{\mathsf T_1} and \NSi{NS-F-P092-037}{\mathsf T_2} will also recover radial
fluxes of axial and azimuthal momentum. Write \NSi{NS-F-P092-038}{\Psi} for the azimuthal
coefficient of the vector potential. The physical normalizations, with a star denoting a normalized
representative, are
\NSFormula{NS-F-P092-039}{display}{8.11}
\begin{equation}
\begin{gathered}
u_*=Q^A u_{\mathrm{phys}},\qquad
g_*=Q^{2A+1/2}g_{\mathrm{phys}},\qquad
p_*=Q^{2A}p_{\mathrm{phys}},\\
\Psi_*=Q^{A-1/2}\Psi_{\mathrm{phys}},\qquad
\rho_*=Q^{1/2}\rho_{\mathrm{phys}}.
\end{gathered}
\tag{8.11}\label{eq:8.11}
\end{equation}
In particular,
\NSi{NS-F-P092-040}{Q^{2A}\mathcal I_{c,\mathrm{phys}}g_{\mathrm{phys}}=\mathcal I_cg_*}.
Preservation of the \NSi{NS-F-P092-041}{\varepsilon} exponent refers to these normalized
representatives, with the integration length included.

\begin{proposition}\label{prop:8.3}
Fix the radial operators of Lemma~\ref{lem:8.2} and the bump \NSi{NS-F-P092-042}{\rho} above.
\NSPage{093}%
\begin{enumerate}[label=(\roman*)]
\item \emph{Pressure.} For a smooth shell-supported scalar source
\NSi{NS-F-P093-001}{g_r(R,Z,T,y)}, define the radial defect \NSi{NS-F-P093-002}{\mathcal P}
and pressure \NSi{NS-F-P093-003}{p_m} by
\NSFormula{NS-F-P093-004}{display}{8.12}
\begin{equation}
\mathcal P=\int\langle g_r\rangle_Y\,dR,
\qquad p_m=\mathsf T_0(g_r-\rho\mathcal P),
\qquad \int\rho\,dR=1.
\tag{8.12}\label{eq:8.12}
\end{equation}
Then
\NSFormula{NS-F-P093-005}{display}{8.13}
\begin{equation}
D_rp_m-g_r=-\rho\mathcal P-\mathcal A_0(g_r-\rho\mathcal P),
\qquad \partial_R\langle p_m\rangle_Y=\langle g_r\rangle_Y-\rho\mathcal P.
\tag{8.13}\label{eq:8.13}
\end{equation}
If \NSi{NS-F-P093-006}{g_r\in \mathsf M_\alpha}, then
\NSi{NS-F-P093-007}{\mathcal P\in \mathsf S_\alpha} and
\NSi{NS-F-P093-008}{p_m\in \mathsf M_\alpha}. A change in \NSi{NS-F-P093-009}{g_r} of class
\NSi{NS-F-P093-010}{\mathsf M_\alpha} produces a pressure change of that same class. Under this
class hypothesis, the cutoff remainder in \eqref{eq:8.13} has zero auxiliary mean and is flat as
\NSi{NS-F-P093-011}{q\downarrow0} at every fixed derivative order.

\item \emph{Axial velocity.} Let \NSi{NS-F-P093-012}{\gamma_d(R,Z,T,y)} be a smooth
shell-supported desired axial increment satisfying
\NSi{NS-F-P093-013}{\int R\langle\gamma_d\rangle_Y\,dR=0} at every
\NSi{NS-F-P093-014}{(Z,T)}. Define the physical azimuthal vector-potential coefficient
\NSi{NS-F-P093-015}{\Psi=-r^{-1}\mathcal I_c(r\gamma_{d,\mathrm{phys}})}, using \eqref{eq:8.4},
and define its normalized velocity increments by
\NSFormula{NS-F-P093-016}{display}{8.14}
\begin{equation}
\Psi_*=\mathsf T_1\gamma_d,
\qquad \Delta\beta=-\varepsilon\partial_Z\Psi_*,
\qquad \Delta\gamma=(D_r+1/R)\Psi_*=\gamma_d-\mathcal A_1\gamma_d.
\tag{8.14}\label{eq:8.14}
\end{equation}
The actual increment \NSi{NS-F-P093-017}{(\Delta\beta,0,\Delta\gamma)} is divergence-free and
has \NSi{NS-F-P093-018}{\langle\Delta\gamma\rangle_Y=\langle\gamma_d\rangle_Y}. If
\NSi{NS-F-P093-019}{\gamma_d\in \mathsf M_\alpha}, then
\NSi{NS-F-P093-020}{\Psi_*,\Delta\gamma\in \mathsf M_\alpha} and
\NSi{NS-F-P093-021}{\Delta\beta\in \mathsf M_{\alpha+1}}. If
\NSi{NS-F-P093-022}{\gamma_d} is independent of \NSi{NS-F-P093-023}{Y}, its zero weighted
integral implies \NSi{NS-F-P093-024}{\Delta\gamma=\gamma_d} pointwise. The pressure and vector
potential are compactly supported in the active shell and retain the projection of the source support
to \NSi{NS-F-P093-025}{(Z,T)}.
\end{enumerate}
\end{proposition}

\begin{proof}
\noindent\textbf{Part (i): pressure.} The source in \eqref{eq:8.12} has zero radial auxiliary-mean
integral:
\NSFormula{NS-F-P093-026}{display}{none}
\[
\int\langle g_r-\rho\mathcal P\rangle_Y\,dR
=\mathcal P-\mathcal P\int\rho\,dR=0.
\]
Apply Lemma~\ref{lem:8.2} with \NSi{NS-F-P093-027}{e=0}, then take its auxiliary average. The
moment \NSi{NS-F-P093-028}{\mathcal P} lies in \NSi{NS-F-P093-029}{\mathsf S_\alpha} when
\NSi{NS-F-P093-030}{g_r\in \mathsf M_\alpha}: integrate its derivatives of any fixed order over
the bounded shell, using the smooth zero extension and the boundedness of
\NSi{NS-F-P093-031}{\zeta\delta^{-B}} for each finite \NSi{NS-F-P093-032}{B}. The interior bump
\NSi{NS-F-P093-033}{\rho\mathcal P} lies in \NSi{NS-F-P093-034}{\mathsf M_\alpha}. Since
\NSi{NS-F-P093-035}{\rho} is fixed independently of the current velocity, the same argument proves
the difference statement and completes part (i).

\noindent\textbf{Part (ii): axial velocity.} The azimuthal potential generates the two physical
velocity components \NSi{NS-F-P093-036}{-\partial_z\Psi} and
\NSi{NS-F-P093-037}{(\mathfrak r+1/r)\Psi}. Commutation of the operators
\NSi{NS-F-P093-038}{\mathfrak r} and \NSi{NS-F-P093-039}{\partial_z}, which differentiate after
evaluation at the phase map, proves that their divergence vanishes exactly:
\NSFormula{NS-F-P093-040}{display}{none}
\[
(\mathfrak r+r^{-1})(-\partial_z\Psi)+\partial_z\bigl((\mathfrak r+r^{-1})\Psi\bigr)=0.
\]
The weighted primitive identity gives the last equality in \eqref{eq:8.14}, including its sign. The
zero weighted axial-flux integral gives zero auxiliary average of its cutoff remainder. Finally
\NSi{NS-F-P093-041}{Q^{1/2-D}=Q^h=\varepsilon}, which gives the extra factor of
\NSi{NS-F-P093-042}{\varepsilon} in \NSi{NS-F-P093-043}{\Delta\beta} in \eqref{eq:8.14}. The
remaining claims follow from Lemma~\ref{lem:8.2}, completing part (ii). No slow derivative is
commuted through \NSi{NS-F-P093-044}{\mathcal I_c} in this argument.
\end{proof}

\subsection{Three compatibility defects in the integrated equations.}\label{sec:8.4}
In this subsection we identify the integral defects left by pressure reconstruction \eqref{eq:8.13}
and the two tangential equations in \eqref{eq:8.3}. We will then construct tangential stresses whose
radial divergences cancel the auxiliary-averaged residuals up to two explicit bump terms
(Corollary~\ref{cor:8.5}). Pressure reconstruction leaves the integral
\NSi{NS-F-P093-045}{\mathcal P} in \eqref{eq:8.12} to be corrected. The auxiliary averages of the
tangential residuals must likewise have zero weighted radial integrals before they can be represented
by compactly supported stresses (Lemma~\ref{lem:8.2}). The moment constraints \eqref{eq:8.2}
remove the time and axial-viscosity terms from those integrals. The surviving axial fluxes give two
further defects, as in Equations~\eqref{eq:8.15} and \eqref{eq:8.16}.

\NSPage{094}%
Alongside \NSi{NS-F-P094-001}{\mathcal P} in \eqref{eq:8.12}, define the two flux defects and the
second radial moment \NSi{NS-F-P094-002}{c_\rho} of \NSi{NS-F-P094-003}{\rho} by
\NSFormula{NS-F-P094-004}{display}{8.15}
\begin{equation}
\begin{aligned}
\mathcal J_\theta
  &=\int R^2\langle Gv+V\gamma+\gamma v+W_{z\theta}\rangle_Y\,dR,\\
\mathcal J_z
  &=\int R\langle2G\gamma+\gamma^2+W_{zz}\rangle_Y\,dR
    -\frac12\int R^2\langle g_r\rangle_Y\,dR,
&c_\rho&=\frac12\int R^2\rho\,dR.
\end{aligned}
\tag{8.15}\label{eq:8.15}
\end{equation}

\begin{proposition}\label{prop:8.4}
Let \NSi{NS-F-P094-005}{(\beta,v,\gamma)} and \NSi{NS-F-P094-006}{w} be the smooth
shell-supported corrections of \eqref{eq:8.1}. Suppose that the two moments in \eqref{eq:8.2}
vanish at every \NSi{NS-F-P094-007}{(Z,T)}, and reconstruct \NSi{NS-F-P094-008}{p_m} from their
radial source \NSi{NS-F-P094-009}{g_r} by \eqref{eq:8.12}. For the residual components
\NSi{NS-F-P094-010}{E_\theta,E_z} in \eqref{eq:8.3} and the defects
\NSi{NS-F-P094-011}{\mathcal J_\theta,\mathcal J_z,c_\rho} in \eqref{eq:8.15}, one has
\NSFormula{NS-F-P094-012}{display}{8.16}
\begin{equation}
\int R^2\langle E_\theta\rangle_Y\,dR
=\varepsilon\partial_Z\mathcal J_\theta,
\qquad
\int R\langle E_z\rangle_Y\,dR
=\varepsilon\partial_Z(\mathcal J_z+c_\rho\mathcal P).
\tag{8.16}\label{eq:8.16}
\end{equation}
\end{proposition}

\begin{proof}
Auxiliary averaging turns \NSi{NS-F-P094-013}{D_r} into \NSi{NS-F-P094-014}{\partial_R} and
\NSi{NS-F-P094-015}{\mathfrak t_*} into \NSi{NS-F-P094-016}{-\varepsilon\partial_T}. Each radial
divergence in \eqref{eq:8.3}, including that of the residual stress tensor, is an endpoint term after
multiplication by its indicated weight. Those terms vanish by compact support. The radial viscous
integrals vanish by two integrations by parts:
\NSFormula{NS-F-P094-017}{display}{none}
\[
\int R^2(v_{RR}+R^{-1}v_R-R^{-2}v)\,dR
=(2-1-1)\int v\,dR=0,
\]
\NSFormula{NS-F-P094-018}{display}{none}
\[
\int R(\gamma_{RR}+R^{-1}\gamma_R)\,dR
=-\int\gamma_R\,dR+\int\gamma_R\,dR=0.
\]
These formulas are applied to the auxiliary averages. The time terms and axial viscosity terms are
derivatives of the two zero integrals. The pressure moment is fixed by its compactly supported
reconstruction:
\NSFormula{NS-F-P094-019}{display}{none}
\[
\int R\langle p_m\rangle_Y\,dR
=-\frac12\int R^2\partial_R\langle p_m\rangle_Y\,dR
=-\frac12\int R^2\langle g_r\rangle_Y\,dR+c_\rho\mathcal P.
\]
Substitution in the surviving axial fluxes yields \eqref{eq:8.16}. Since
\NSi{NS-F-P094-020}{\rho} is scaled by the moving \NSi{NS-F-P094-021}{q},
\NSi{NS-F-P094-022}{c_\rho} can depend on \NSi{NS-F-P094-023}{(Z,T)}. Its full product with
\NSi{NS-F-P094-024}{\mathcal P} therefore remains inside \NSi{NS-F-P094-025}{\partial_Z}.
\end{proof}

The moments in physical variables and their normalized representatives are related by
\NSFormula{NS-F-P094-026}{display}{8.17}
\begin{equation}
\begin{gathered}
(M_\theta)_*=Q^{A-3/2}(M_\theta)_{\mathrm{phys}},
\qquad (M_z)_*=Q^{A-1}(M_z)_{\mathrm{phys}},\\
\mathcal P_*=Q^{2A}\mathcal P_{\mathrm{phys}},
\qquad (c_\rho)_*=Q^{-1}(c_\rho)_{\mathrm{phys}},\\
(\mathcal J_\theta)_*=Q^{2A-3/2}(\mathcal J_\theta)_{\mathrm{phys}},
\qquad (\mathcal J_z)_*=Q^{2A-1}(\mathcal J_z)_{\mathrm{phys}}.
\end{gathered}
\tag{8.17}\label{eq:8.17}
\end{equation}
The zero integral constraints are invariant under normalization. Each nonzero correction target uses
the normalization of its corresponding row.

The two tangential residual moments are therefore axial derivatives of
\NSi{NS-F-P094-027}{\mathcal J_\theta} and
\NSi{NS-F-P094-028}{\mathcal J_z+c_\rho\mathcal P}. We subtract a fixed normalized bump times
each weighted residual integral, then apply the radial inverse. The sources are independent of the
auxiliary variable and have zero weighted integrals, so this inverse is exact. The resulting pair will
be the target for an \emph{added covariance} of the waves: its weighted radial divergence enters
\eqref{eq:8.3} with a positive sign through \NSi{NS-F-P094-029}{W_{r\theta}} and
\NSi{NS-F-P094-030}{W_{rz}}. Fix interior profiles
\NSi{NS-F-P094-031}{\widehat\sigma_\theta,\widehat\sigma_z} with
\NSi{NS-F-P094-032}{\int\xi^2\widehat\sigma_\theta(\xi)\,d\xi
=\int\xi\widehat\sigma_z(\xi)\,d\xi=1}, and define
\NSFormula{NS-F-P094-033}{display}{none}
\[
\sigma_{\theta,\mathrm{phys}}=q^{-3/2}\widehat\sigma_\theta(r/\sqrt q),
\qquad
\sigma_{z,\mathrm{phys}}=q^{-1}\widehat\sigma_z(r/\sqrt q).
\]
Their supports lie in the mean patch. In each chart use
\NSi{NS-F-P094-034}{\sigma_\theta=Q^{3/2}\sigma_{\theta,\mathrm{phys}}} and
\NSi{NS-F-P094-035}{\sigma_z=Q\sigma_{z,\mathrm{phys}}}, so
\NSi{NS-F-P094-036}{\int R^2\sigma_\theta\,dR=\int R\sigma_z\,dR=1}. This fixes the same physical
bump functions on chart overlaps.
\NSPage{095}%
\begin{corollary}\label{cor:8.5}
Under the hypotheses of Proposition~\ref{prop:8.4}, define the covariance targets
\NSFormula{NS-F-P095-001}{display}{8.18}
\begin{equation}
\begin{aligned}
H_\theta&=-\mathsf T_2\bigl(\langle E_\theta\rangle_Y
  -\sigma_\theta\varepsilon\partial_Z\mathcal J_\theta\bigr),\\
H_z&=-\mathsf T_1\bigl(\langle E_z\rangle_Y
  -\sigma_z\varepsilon\partial_Z(\mathcal J_z+c_\rho\mathcal P)\bigr).
\end{aligned}
\tag{8.18}\label{eq:8.18}
\end{equation}
They are independent of the auxiliary variable, compactly supported in the active shell, and satisfy
\NSFormula{NS-F-P095-002}{display}{none}
\begin{align*}
(D_r+2/R)H_\theta&=-\langle E_\theta\rangle_Y
  +\sigma_\theta\varepsilon\partial_Z\mathcal J_\theta,\\
(D_r+1/R)H_z&=-\langle E_z\rangle_Y
  +\sigma_z\varepsilon\partial_Z(\mathcal J_z+c_\rho\mathcal P).
\end{align*}
The radial primitives act at fixed \NSi{NS-F-P095-003}{(Z,T)}. Thus
\NSi{NS-F-P095-004}{H_\theta} and \NSi{NS-F-P095-005}{H_z} vanish whenever their respective
integrands in \eqref{eq:8.18} vanish for every \NSi{NS-F-P095-006}{R}; the construction does not
enlarge support in the axial variable \NSi{NS-F-P095-007}{Z}.
\end{corollary}

\begin{proof}
The integrated identities and the normalization of the two bumps give
\NSFormula{NS-F-P095-008}{display}{none}
\begin{align*}
\int R^2\bigl(\langle E_\theta\rangle_Y
  -\sigma_\theta\varepsilon\partial_Z\mathcal J_\theta\bigr)\,dR&=0,\\
\int R\bigl(\langle E_z\rangle_Y
  -\sigma_z\varepsilon\partial_Z(\mathcal J_z+c_\rho\mathcal P)\bigr)\,dR&=0.
\end{align*}
Both sources are independent of the auxiliary variable. Their cutoff remainders vanish by
Lemma~\ref{lem:8.2}, which gives the two identities and the support assertions.
\end{proof}

These radial divergence identities retain the factors
\NSi{NS-F-P095-009}{\varepsilon\partial_Z\mathcal J_\theta} and
\NSi{NS-F-P095-010}{\varepsilon\partial_Z(\mathcal J_z+c_\rho\mathcal P)} in the remaining bump
terms. We can therefore correct the three defects at fixed \NSi{NS-F-P095-011}{(Z,T)} without
losing the displayed factors of \NSi{NS-F-P095-012}{\varepsilon\partial_Z}.

\subsection{Inverting the auxiliary-time derivative on zero-mean functions.}\label{sec:8.5}
Corollary~\ref{cor:8.5} reduces the auxiliary-averaged correction to the three integral defects and a
compactly supported stress. The complementary part
\NSi{NS-F-P095-013}{(E_\theta^\circ,E_z^\circ)} of the mean residual has zero auxiliary average at
every slow point, by \eqref{eq:3.9}. In this subsection we construct its correction through the fast
time derivative, whose Fourier multiplier is invertible on nonzero torus frequencies
(Equation~\eqref{eq:6.7}). The axial velocity increment is then realized through its azimuthal
potential in \eqref{eq:8.14} to preserve incompressibility.

Recall from Equations~\eqref{eq:6.4} and \eqref{eq:6.6} that
\NSi{NS-F-P095-014}{N_{\mathrm{abs}}=v_t\cdot\partial_Y} acts in the original auxiliary coordinate.
On a common torus \NSi{NS-F-P095-015}{y=Y_{i_0}}, the normalized fast time derivative is
\NSi{NS-F-P095-016}{c_{i_0}N_{i_0}}, where
\NSi{NS-F-P095-017}{N_{i_0}=v_t\cdot\partial_y} and
\NSi{NS-F-P095-018}{c_{i_0}=T_g^{i_0}Q^{1+h}}. Here the normalized physical time derivative splits
as \NSi{NS-F-P095-019}{\mathfrak t_*=-\varepsilon\partial_T+c_{i_0}N_{i_0}}. The slow term
\NSi{NS-F-P095-020}{-\varepsilon\partial_T} differentiates the coefficient's explicit time
dependence while holding \NSi{NS-F-P095-021}{y} fixed. The fast term
\NSi{NS-F-P095-022}{c_{i_0}N_{i_0}} differentiates its auxiliary oscillations along the physical
phase map. We invert the fast term here and retain the slow term in the residual.

\begin{lemma}\label{lem:8.6}
Let \NSi{NS-F-P095-023}{N=v_t\cdot\partial_y} on \NSi{NS-F-P095-024}{\mathbb T^2}, with
\NSi{NS-F-P095-025}{v_t} as in \eqref{eq:6.2}, and let \NSi{NS-F-P095-026}{F} be a smooth function
of zero Haar mean. Then \NSi{NS-F-P095-027}{N\varphi=F} has a unique smooth zero-mean solution
\NSi{NS-F-P095-028}{\varphi=N^{-1}F}. This solution is given by the following Fourier series and
satisfies the displayed norm bound for every integer \NSi{NS-F-P095-029}{m\ge0}:
\NSFormula{NS-F-P095-030}{display}{8.19}
\begin{equation}
N^{-1}F=\sum_{k\in\mathbb Z^2\setminus\{0\}}
\frac{\widehat F(k)}{2\pi i v_t\cdot k}e^{2\pi i k\cdot y},
\qquad
\|N^{-1}F\|_{C_y^m}\le C_m\|F\|_{C_y^{m+4}}.
\tag{8.19}\label{eq:8.19}
\end{equation}
For a pair \NSi{NS-F-P095-031}{E_\theta,E_z\in \mathsf M_\alpha} of normalized angularly invariant
residual coefficients, use this inverse in the original auxiliary coordinate
\NSi{NS-F-P095-032}{Y} to define the physical tangential update
\NSi{NS-F-P095-033}{-N_{\mathrm{abs}}^{-1}
  (E_{\theta,\mathrm{phys}}^\circ,E_{z,\mathrm{phys}}^\circ)}. Here
\NSPage{096}%
residuals and velocities have the normalizations in \eqref{eq:8.11}. This physical definition ensures
agreement on chart overlaps. Its chart representatives are
\NSFormula{NS-F-P096-001}{display}{8.20}
\begin{equation}
\Delta v=-c_{i_0}^{-1}N_{i_0}^{-1}E_\theta^\circ,
\qquad
\gamma_d=-c_{i_0}^{-1}N_{i_0}^{-1}E_z^\circ,
\qquad
c_{i_0}^{-1}\le CS_*.
\tag{8.20}\label{eq:8.20}
\end{equation}
Both desired increments \NSi{NS-F-P096-002}{\Delta v,\gamma_d} lie in
\NSi{NS-F-P096-003}{\mathsf M_\alpha} and have zero auxiliary mean at each slow and radial point.
The realization of \NSi{NS-F-P096-004}{\gamma_d} by \eqref{eq:8.14} gives the actual velocity
increment
\NSFormula{NS-F-P096-005}{display}{8.21}
\begin{equation}
\begin{gathered}
(\Delta\beta,\Delta v,\Delta\gamma)
=(-\varepsilon\partial_Z\mathsf T_1\gamma_d,\Delta v,\gamma_d-\mathcal A_1\gamma_d),\\
(\Delta\beta,\Delta v,\Delta\gamma)
\in \mathsf M_{\alpha+1}\times \mathsf M_\alpha\times \mathsf M_\alpha.
\end{gathered}
\tag{8.21}\label{eq:8.21}
\end{equation}
This field is divergence-free and its addition preserves both integral constraints in \eqref{eq:8.2}.
If \NSi{NS-F-P096-006}{a=-\mathcal A_1\gamma_d} denotes the cutoff remainder in the axial
increment, then its fast-time cancellation equations are
\NSFormula{NS-F-P096-007}{display}{8.22}
\begin{equation}
c_{i_0}N_{i_0}\Delta v=-E_\theta^\circ,
\qquad
c_{i_0}N_{i_0}\Delta\gamma=-E_z^\circ+c_{i_0}N_{i_0}a.
\tag{8.22}\label{eq:8.22}
\end{equation}
The angular cancellation is exact, and \NSi{NS-F-P096-008}{c_{i_0}N_{i_0}a} is flat at every fixed
coefficient derivative order.
\end{lemma}

\begin{proof}
The divisor bound \eqref{eq:6.7} gives
\NSi{NS-F-P096-009}{|v_t\cdot k|^{-1}\le C(1+|k|)}. Four more derivatives of
\NSi{NS-F-P096-010}{F} than the requested output leave a summable
\NSi{NS-F-P096-011}{(1+|k|)^{-3}} bound on the Fourier series in two dimensions. This proves
\eqref{eq:8.19}, smoothness, and uniqueness after fixing the zero Fourier coefficient. The multiplier
commutes with every slow coefficient derivative and preserves reality. It acts at fixed
\NSi{NS-F-P096-012}{(R,Z,T)}, so it preserves slow and radial support and their smooth zero
extensions, although it need not preserve auxiliary-torus support.

Let \NSi{NS-F-P096-013}{P_iF(Y)=F(J_g^iY)}. A mode \NSi{NS-F-P096-014}{k} on the
\NSi{NS-F-P096-015}{i}-torus has frequency \NSi{NS-F-P096-016}{(J_g^i)^Tk} in the original
auxiliary coordinate \NSi{NS-F-P096-017}{Y}. Since
\NSi{NS-F-P096-018}{J_gv_t=T_gv_t},
\NSFormula{NS-F-P096-019}{display}{8.23}
\begin{equation}
N_{\mathrm{abs}}P_i=T_g^iP_iN_i,
\qquad
N_{\mathrm{abs}}^{-1}P_i=T_g^{-i}P_iN_i^{-1}.
\tag{8.23}\label{eq:8.23}
\end{equation}
The finite covering preserves Haar average, as follows also from its injective map of frequency
lattices. Applying \eqref{eq:8.19} in the common frequency coordinate
\NSi{NS-F-P096-020}{k} gives the same fixed derivative loss in every chart.

In physical variables the update is
\NSi{NS-F-P096-021}{-N_{\mathrm{abs}}^{-1}E_{\mathrm{phys}}^\circ}. Its chart velocity value has
factor
\NSi{NS-F-P096-022}{Q^{A-(2A+1/2)}T_g^{-i_0}=Q^{-1-h}T_g^{-i_0}=c_{i_0}^{-1}}, which proves
\eqref{eq:8.20}. This factor is bounded by \NSi{NS-F-P096-023}{CS_*} by the choice of covering
level and its bounded difference from a band level. For each output derivative, four additional
derivatives in the torus variables suffice, so the class estimates at every fixed derivative order
retain their \NSi{NS-F-P096-024}{\varepsilon} exponent. Applying the slow-time term
\NSi{NS-F-P096-025}{-\varepsilon\partial_T} to these increments gives an additional factor of
\NSi{NS-F-P096-026}{\varepsilon}.

The desired increments have zero auxiliary mean at each radial point. Hence both desired integral
constraints hold, and the azimuthal-potential realization preserves them exactly. The bounds for
\NSi{NS-F-P096-027}{\Delta\beta} and \NSi{NS-F-P096-028}{\Delta\gamma} in \eqref{eq:8.21}
follow from Proposition~\ref{prop:8.3}(ii). The fast derivative commutes with
\NSi{NS-F-P096-029}{\mathcal I}, \NSi{NS-F-P096-030}{\mathcal J}, and
\NSi{NS-F-P096-031}{\mathcal I_c}: the shifts are translations and
\NSi{NS-F-P096-032}{\chi_m} is torus independent. Applying it to
\NSi{NS-F-P096-033}{\Delta\gamma=\gamma_d+a} proves \eqref{eq:8.22}. Its last term is flat in every
fixed coefficient derivative by \eqref{eq:8.8}. Full slow differentiation need not commute with
\NSi{NS-F-P096-034}{\mathcal I_c}; those derivatives are estimated by Lemma~\ref{lem:8.2} instead.
\end{proof}

\subsection{Five-dimensional correction of the moment constraints.}\label{sec:8.6}
We have constructed inverses for pressure (Proposition~\ref{prop:8.3}), the auxiliary-averaged
tangential residual (Corollary~\ref{cor:8.5}), and the part with zero auxiliary average
(Lemma~\ref{lem:8.6}). The remaining quantities are the three integral defects
\NSi{NS-F-P096-035}{\mathcal P,\mathcal J_\theta,\mathcal J_z} in Equations~\eqref{eq:8.12} and
\eqref{eq:8.15}. In this subsection we will correct their linear changes by a velocity increment that
also preserves the two constraints \eqref{eq:8.2}. Three azimuthal coefficients and two axial
coefficients will solve the five equations \eqref{eq:8.25}.
